\scanpage{160}
Dans le cas favorable envisagé à la fin du n° 3, ($\unsure{J|T}$ de car. nulle, $\Omega_{S/T}$ plat), on peut de $=$ exprimer les $\mu_D$ en termes des $\unsure{\nu|v}_\Delta$. Pour ceci, il suffit de voir que, pour $S$ affine, chaque $D$ provient bien d'un $\Delta$. Ceci résulte du fait suivant (lui-$=$ cas particulier d'un fait bien plus \unsure{générale|général} sur les op. diff) : \deleted{Tout op.}

\uline{Lemme}. Sous les conditions précédentes ($S$ car. nulle, $\Omega_{S/T}$ plat), \deleted{tout} \inserted{les} op diff $E\longrightarrow F$ qui sont tels que $P^\infty_{S/T}(E)\longrightarrow F$ s'annule sur $\delta(E)$, sont exactement ceux \deleted{qui \unsure{se}} de la forme $\Delta\circ d^0_E$, avec $\Delta:E\otimes\Omega_{S/T}\mathrel{\text{---}}F$ un op différentiel. L'indétermination est \unsure{exactement|exactem\textsuperscript{t}} \deleted{\unsure{l'op}} par un op diff. \deleted{\unsure{$D', E\otimes\Omega_{S/T}^2$|$\Delta'\circ d^1_E:E\otimes\Omega_{S/T}$}} \inserted{$\Delta'\circ d^1_E$}, où $\Delta':E\otimes\Omega_{S/T}\longrightarrow F$.

\deleted{Ceci résulte du fait que \unsure{$E\longrightarrow\Omega E$|$E\longrightarrow\Omega$}.}
% archive sheet number: 159
